
This document is based on \textit{Datasheets for Datasets} by~\citep{gebru2021datasheets}.

\subsection{Motivation}

    \textcolor{\sectioncolor}{\textbf{For what purpose was the dataset created?
    }
    Was there a specific task in mind? Was there
    a specific gap that needed to be filled? Please provide a description.
    } \\
    %%%
The dataset was created to support user-centered semantic search in Lean by capturing the diverse intents of mathematicians. It includes synthesized user queries grounded in formal statements, informalized statements from various libraries and research-linked repositories, and augmented proof states. The purpose is to create retrieval models that better reflect real user needs and integrate seamlessly with LLM-based theorem provers. The dataset bridges real user queries and raw informalizations for search engines, fostering research in retrieval-augmented theorem proving and intent-aware formal reasoning.
    % (see \cref{sec:charxiv_overview})
    \\
    %%% 
    
    \textcolor{\sectioncolor}{\textbf{Who created this dataset (e.g., which team, research group) and on behalf
    of which entity (e.g., company, institution, organization)?
    }
    } \\
    %%%
    The dataset was created by the anonymous authors of this Lean Finder project, affiliated with a research group focused on AI-for-math applications. \\
    %%% 
    
    \textcolor{\sectioncolor}{\textbf{What support was needed to make this dataset?
    }
    (e.g.who funded the creation of the dataset? If there is an associated
    grant, provide the name of the grantor and the grant name and number, or if
    it was supported by a company or government agency, give those details.)
    } \\
    %%%
The creation of the dataset was supported by research funding for developing novel applications of LLMs in formal mathematics. Additional support included access to computational resources for fine-tuning and coverage of API costs for data synthesis and processing. \\
    %%% 
    
    \textcolor{\sectioncolor}{\textbf{Any other comments?
    }} \\
    %%%
    N/A \\
    %%%

%%%%%%%%%%%%%%%%%%%%%%%%%%%%%%%%%%%%%%%%%%%%%%%%%%%%%%%%%%%%%%%%%%%%%%%%%%%%%%%%
\subsection{Composition}
    \textcolor{\sectioncolor}{\textbf{What do the instances that comprise the dataset represent (e.g., documents,
    photos, people, countries)?
    }
    Are there multiple types of instances (e.g., movies, users, and ratings;
    people and interactions between them; nodes and edges)? Please provide a
    description.
    } \\
    %%%
Each instance represents a Lean-related query, which includes statement pair or proof context designed to capture mathematicians’ intents. The dataset includes multiple types of instances: (i) synthesized user queries grounded in formal Lean statements, (ii) informalized statements providing natural language descriptions of Lean theorems, (iii) augmented proof states representing transitions between proof steps, and (iv) formal statement extracted from Lean libraries and research-linked repositories. Together, these instance types reflect different ways mathematicians interact with Lean and enable comprehensive training for semantic retrieval models.

    %%% 
    
    \textcolor{\sectioncolor}{\textbf{How many instances are there in total (of each type, if appropriate)?
    }
    } \\
    %%%
The dataset comprises over 1.4 million query–code pairs. Specifically, it includes 582k synthesized user queries, 244k informalized statements, 338k augmented proof states, and 244k formal statement.

    % More details are shown in \cref{tab:overall_statistics}.

    %%% 
    \textcolor{\sectioncolor}{\textbf{Does the dataset contain all possible instances or is it a sample (not
    necessarily random) of instances from a larger set?
    }
    If the dataset is a sample, then what is the larger set? Is the sample
    representative of the larger set (e.g., geographic coverage)? If so, please
    describe how this representativeness was validated/verified. If it is not
    representative of the larger set, please describe why not (e.g., to cover a
    more diverse range of instances, because instances were withheld or
    unavailable).
    } \\
    %%%
    The dataset is a curated sample rather than the complete set of possible Lean queries. The larger set consists of all formal statements and proofs in mathlib4 and research-linked Lean repositories, along with the open-ended space of user queries posed by mathematicians. Since real user queries are limited in number and often lack ground-truth resolutions, we synthesize queries grounded in formal statements to approximate this space. While not exhaustive, the dataset is designed to be broadly representative by covering multiple input modalities (synthetic user queries, informalizations, proof states, and formal statements) and by drawing from both mathlib and research-linked repositories to capture diverse mathematical content.
 \\
    %%%
    
    \textcolor{\sectioncolor}{\textbf{What data does each instance consist of?
    }
    “Raw” data (e.g., unprocessed text or images) or features? In either case,
    please provide a description.
    } \\
    %%%
Each instance consists of raw text in the form of Lean 4 formal statements, their associated natural language variants, or proof contexts. Depending on modality, this may include: (i) synthesized user queries phrased in natural language, (ii) informalized statements produced from formal Lean code, (iii) augmented proof states describing proof progressions, and (iv) formal statement. All data are stored as text pairs.

    %%% 
    
    \textcolor{\sectioncolor}{\textbf{Is there a label or target associated with each instance?
    }
    If so, please provide a description.
    } \\
    %%%
    Yes. Each query is paired with one or more target Lean 4 formal statements that serve as the ground-truth retrieval target. In the case of augmented proof states, the target can be one or more Lean statements used in the corresponding proof transition. For informalizations, synthetic user queries, formal statement, the target is the complete Lean statement from which they were derived.
 \\
    %%% 
    
    \textcolor{\sectioncolor}{\textbf{Is any information missing from individual instances?
    }
    If so, please provide a description, explaining why this information is
    missing (e.g., because it was unavailable). This does not include
    intentionally removed information, but might include, e.g., redacted text.
    } \\
    %%%
    N/A \\
    %%% 
    
    \textcolor{\sectioncolor}{\textbf{Are relationships between individual instances made explicit (e.g., users’
    movie ratings, social network links)?
    }
    If so, please describe how these relationships are made explicit.
    } \\
    %%%
    No. Each instance is an independent query–code pair. \\
    %%% 
    
    \textcolor{\sectioncolor}{\textbf{Are there recommended data splits (e.g., training, development/validation,
    testing)?
    }
    If so, please provide a description of these splits, explaining the
    rationale behind them.
    } \\
    %%%
    Yes, the dataset is split into training and testing sets. \\
    %%% 
    
    \textcolor{\sectioncolor}{\textbf{Are there any errors, sources of noise, or redundancies in the dataset?
    }
    If so, please provide a description.
    } \\
    %%%
    Yes. The synthesized queries from formal statements are generated using LLMs, some may contain minor phrasing inconsistencies or unnatural wording compared to real user queries. Informalizations may also result in inaccuracies.
 \\
    %%% 
    
    \textcolor{\sectioncolor}{\textbf{Is the dataset self-contained, or does it link to or otherwise rely on
    external resources (e.g., websites, tweets, other datasets)?
    }
    If it links to or relies on external resources, a) are there guarantees
    that they will exist, and remain constant, over time; b) are there official
    archival versions of the complete dataset (i.e., including the external
    resources as they existed at the time the dataset was created); c) are
    there any restrictions (e.g., licenses, fees) associated with any of the
    external resources that might apply to a future user? Please provide
    descriptions of all external resources and any restrictions associated with
    them, as well as links or other access points, as appropriate.
    } \\
    %%%
    The dataset is self-contained, with no reliance on external or dynamic resources. \\
    %%% 
    
    \textcolor{\sectioncolor}{\textbf{Does the dataset contain data that might be considered confidential (e.g.,
    data that is protected by legal privilege or by doctor-patient
    confidentiality, data that includes the content of individuals’ non-public
    communications)?
    }
    If so, please provide a description.
    } \\
    %%%
    No. We do not release any real user queries that we collected for fine-tuning and data synthesis. \\
    %%% 
    
    \textcolor{\sectioncolor}{\textbf{Does the dataset contain data that, if viewed directly, might be offensive,
    insulting, threatening, or might otherwise cause anxiety?
    }
    If so, please describe why.
    } \\
    %%%
    No. \\
    %%% 
    
    \textcolor{\sectioncolor}{\textbf{Does the dataset relate to people?
    }
    If not, you may skip the remaining questions in this section.
    } \\
    %%%
    No. \\
    %%% 
    
    \textcolor{\sectioncolor}{\textbf{Does the dataset identify any subpopulations (e.g., by age, gender)?
    }
    If so, please describe how these subpopulations are identified and
    provide a description of their respective distributions within the dataset.
    } \\
    %%%
    No. \\
    %%% 
    
    \textcolor{\sectioncolor}{\textbf{Is it possible to identify individuals (i.e., one or more natural persons),
    either directly or indirectly (i.e., in combination with other data) from
    the dataset?
    }
    If so, please describe how.
    } \\
    %%%
    No. \\
    %%% 
    
    \textcolor{\sectioncolor}{\textbf{Does the dataset contain data that might be considered sensitive in any way
    (e.g., data that reveals racial or ethnic origins, sexual orientations,
    religious beliefs, political opinions or union memberships, or locations;
    financial or health data; biometric or genetic data; forms of government
    identification, such as social security numbers; criminal history)?
    }
    If so, please provide a description.
    } \\
    %%%
    No. \\
    %%% 
    
    \textcolor{\sectioncolor}{\textbf{Any other comments?
    }} \\
    %%%
    The dataset’s richness in diversity makes it a significant resource for advancing formal mathematics via AI. \\
    %%%

%%%%%%%%%%%%%%%%%%%%%%%%%%%%%%%%%%%%%%%%%%%%%%%%%%%%%%%%%%%%%%%%%%%%%%%%%%%%%%%%
\subsection{Collection}

    \textcolor{\sectioncolor}{\textbf{How was the data associated with each instance acquired?
    }
    Was the data directly observable (e.g., raw text, movie ratings),
    reported by subjects (e.g., survey responses), or indirectly
    inferred/derived from other data (e.g., part-of-speech tags, model-based
    guesses for age or language)? If data was reported by subjects or
    indirectly inferred/derived from other data, was the data
    validated/verified? If so, please describe how.
    } \\
    %%%
The data was primarily derived from existing Lean 4 formal statements and proofs in mathlib and repositories. From these Lean statements, additional modalities were generated: informalizations, synthetic user queries and augmented proof states were produced by prompting LLMs.\\
    %%% 
    
    \textcolor{\sectioncolor}{\textbf{Over what timeframe was the data collected?
    }
    Does this timeframe match the creation timeframe of the data associated
    with the instances (e.g., recent crawl of old news articles)? If not,
    please describe the timeframe in which the data associated with the
    instances was created. Finally, list when the dataset was first published.
    } \\
    %%%
    All the data is collected in 2025. \\
    %%% 
    
    \textcolor{\sectioncolor}{\textbf{What mechanisms or procedures were used to collect the data (e.g., hardware
    apparatus or sensor, manual human curation, software program, software
    API)?
    }
    How were these mechanisms or procedures validated?
    } \\
    %%%
    We extract Lean statements from mathlib and other Lean repositories, and use GPT-4o API to generate the queries. \\
    %%% 
    
    \textcolor{\sectioncolor}{\textbf{What was the resource cost of collecting the data?
    }
    (e.g. what were the required computational resources, and the associated
    financial costs, and energy consumption - estimate the carbon footprint.)} \\
    %%%
    The total cost for GPT-4o API is around \$2000.
    % gift cards, gpt API
    \\
    %%% 
    
    \textcolor{\sectioncolor}{\textbf{If the dataset is a sample from a larger set, what was the sampling
    strategy (e.g., deterministic, probabilistic with specific sampling
    probabilities)?
    }
    } \\
    %%%
    N/A \\
    %%% 
    
    \textcolor{\sectioncolor}{\textbf{Who was involved in the data collection process (e.g., students,
    crowdworkers, contractors) and how were they compensated (e.g., how much
    were crowdworkers paid)?
    }
    } \\
    %%%
    Graduate students and researchers with expertise in formal mathematics and AI. \\
    %%% 
    
    \textcolor{\sectioncolor}{\textbf{Were any ethical review processes conducted (e.g., by an institutional
    review board)?
    }
    If so, please provide a description of these review processes, including
    the outcomes, as well as a link or other access point to any supporting
    documentation.
    } \\
    %%%
    No. \\
    %%% 
    
    \textcolor{\sectioncolor}{\textbf{Does the dataset relate to people?
    }
    If not, you may skip the remainder of the questions in this section.
    } \\
    %%%
    No. \\
    %%% 
    
    \textcolor{\sectioncolor}{\textbf{Did you collect the data from the individuals in question directly, or
    obtain it via third parties or other sources (e.g., websites)?
    }
    } \\
    %%%
    The Lean statements from our data was collected through mathlib and other Lean related github repositories.\\
    %%% 
    
    \textcolor{\sectioncolor}{\textbf{Were the individuals in question notified about the data collection?
    }
    If so, please describe (or show with screenshots or other information) how
    notice was provided, and provide a link or other access point to, or
    otherwise reproduce, the exact language of the notification itself.
    } \\
    %%%
    N/A \\
    %%% 
    
    \textcolor{\sectioncolor}{\textbf{Did the individuals in question consent to the collection and use of their
    data?
    }
    If so, please describe (or show with screenshots or other information) how
    consent was requested and provided, and provide a link or other access
    point to, or otherwise reproduce, the exact language to which the
    individuals consented.
    } \\
    %%%
    N/A \\
    %%% 
    
    \textcolor{\sectioncolor}{\textbf{If consent was obtained, were the consenting individuals provided with a
    mechanism to revoke their consent in the future or for certain uses?
    }
     If so, please provide a description, as well as a link or other access
     point to the mechanism (if appropriate)
    } \\
    %%%
    N/A \\
    %%% 
    
    \textcolor{\sectioncolor}{\textbf{Has an analysis of the potential impact of the dataset and its use on data
    subjects (e.g., a data protection impact analysis)been conducted?
    }
    If so, please provide a description of this analysis, including the
    outcomes, as well as a link or other access point to any supporting
    documentation.
    } \\
    %%%
    No.\\
    %%% 
    
    \textcolor{\sectioncolor}{\textbf{Any other comments?
    }}\\ N/A \\
    %%%
     % \\
    %%%

%%%%%%%%%%%%%%%%%%%%%%%%%%%%%%%%%%%%%%%%%%%%%%%%%%%%%%%%%%%%%%%%%%%%%%%%%%%%%%%%
\subsection{Preprocessing / Cleaning / Labeling}

    \textcolor{\sectioncolor}{\textbf{Was any preprocessing/cleaning/labeling of the data
    done(e.g.,discretization or bucketing, tokenization, part-of-speech
    tagging, SIFT feature extraction, removal of instances, processing of
    missing values)?
    }
    If so, please provide a description. If not, you may skip the remainder of
    the questions in this section.
    } \\
    %%%
    Yes. Formal Lean statements were extracted from mathlib and related repositories and filtering was applied to keep only the relevant Lean statements. Filtering was also applied to clusters before generating synthetic user queries to avoid invalid synthetic user queries.
 \\
    %%%

    \textcolor{\sectioncolor}{\textbf{Was the “raw” data saved in addition to the preprocessed/cleaned/labeled
    data (e.g., to support unanticipated future uses)?
    }
    If so, please provide a link or other access point to the “raw” data.
    } \\
    %%%
    Yes, raw data and intermediate representations are retained for reproducibility and future use. \\
    %%%

    \textcolor{\sectioncolor}{\textbf{Is the software used to preprocess/clean/label the instances available?
    }
    If so, please provide a link or other access point.
    } \\
    %%%
    The tools and scripts for preprocessing will be released upon acceptance. \\
    %%%

    \textcolor{\sectioncolor}{\textbf{Any other comments?
    }} \\
    %%%
    No. \\
    %%%

%%%%%%%%%%%%%%%%%%%%%%%%%%%%%%%%%%%%%%%%%%%%%%%%%%%%%%%%%%%%%%%%%%%%%%%%%%%%%%%%
\subsection{Uses}

    \textcolor{\sectioncolor}{\textbf{Has the dataset been used for any tasks already?
    }
    If so, please provide a description.
    } \\
    %%%
    Yes, it was used to train and evaluate the Lean Finder and benchmark its performance against other retrieval models. \\
    %%%

    \textcolor{\sectioncolor}{\textbf{Is there a repository that links to any or all papers or systems that use the dataset?
    }
    If so, please provide a link or other access point.
    } \\
    %%%
    N/A \\
    %%%

    \textcolor{\sectioncolor}{\textbf{What (other) tasks could the dataset be used for?
    }
    } \\
    %%%
    Beyond semantic search and retrieval, the dataset could support tasks such as query generation, informal-to-formal translation, retrieval-augmented theorem proving, and exploring human–AI interaction in formal reasoning systems. \\
    %%%

    \textcolor{\sectioncolor}{\textbf{Is there anything about the composition of the dataset or the way it was
    collected and preprocessed/cleaned/labeled that might impact future uses?
    }
    For example, is there anything that a future user might need to know to
    avoid uses that could result in unfair treatment of individuals or groups
    (e.g., stereotyping, quality of service issues) or other undesirable harms
    (e.g., financial harms, legal risks) If so, please provide a description.
    Is there anything a future user could do to mitigate these undesirable
    harms?
    } \\
    %%%
    N/A \\
    %%%

    \textcolor{\sectioncolor}{\textbf{Are there tasks for which the dataset should not be used?
    }
    If so, please provide a description.
    } \\
    %%%
    The dataset is not suitable for tasks unrelated to Lean 4. \\
    %%%

    \textcolor{\sectioncolor}{\textbf{Any other comments?
    }} \\
    %%%
    No. \\
    %%%

%%%%%%%%%%%%%%%%%%%%%%%%%%%%%%%%%%%%%%%%%%%%%%%%%%%%%%%%%%%%%%%%%%%%%%%%%%%%%%%%
\subsection{Distribution}

    \textcolor{\sectioncolor}{\textbf{Will the dataset be distributed to third parties outside of the entity
    (e.g., company, institution, organization) on behalf of which the dataset
    was created?
    }
    If so, please provide a description.
    } \\
    %%%
    Yes, the dataset will be made publicly available for research purposes. \\
    %%%

    \textcolor{\sectioncolor}{\textbf{How will the dataset will be distributed (e.g., tarball on website, API,
    GitHub)?
    }
    Does the dataset have a digital object identifier (DOI)?
    } \\
    %%%
    The dataset will be distributed via GitHub and academic repositories, with accompanying documentation. \\
    %%%

    \textcolor{\sectioncolor}{\textbf{When will the dataset be distributed?
    }
    } \\
    %%%
    The dataset is expected to be released following the ICLR 2026 conference. \\
    %%%

    \textcolor{\sectioncolor}{\textbf{Will the dataset be distributed under a copyright or other intellectual
    property (IP) license, and/or under applicable terms of use (ToU)?
    }
    If so, please describe this license and/or ToU, and provide a link or other
    access point to, or otherwise reproduce, any relevant licensing terms or
    ToU, as well as any fees associated with these restrictions.
    } \\
    %%%
    Yes, it will be distributed under a permissive license (e.g., CC BY-SA 4.0) to encourage research use. \\
    %%%

    \textcolor{\sectioncolor}{\textbf{Have any third parties imposed IP-based or other restrictions on the data
    associated with the instances?
    }
    If so, please describe these restrictions, and provide a link or other
    access point to, or otherwise reproduce, any relevant licensing terms, as
    well as any fees associated with these restrictions.
    } \\
    %%%
    All charts are subjected to their respective copyrights by the authors of this paper. \\
    %%%

    \textcolor{\sectioncolor}{\textbf{Do any export controls or other regulatory restrictions apply to the
    dataset or to individual instances?
    }
    If so, please describe these restrictions, and provide a link or other
    access point to, or otherwise reproduce, any supporting documentation.
    } \\
    %%%
    N/A \\
    %%%

    \textcolor{\sectioncolor}{\textbf{Any other comments?
    }} \\
    %%%
    Distribution will include detailed usage guidelines to ensure proper application of the dataset. \\
    %%%

%%%%%%%%%%%%%%%%%%%%%%%%%%%%%%%%%%%%%%%%%%%%%%%%%%%%%%%%%%%%%%%%%%%%%%%%%%%%%%%%
\subsection{Maintenance}

    \textcolor{\sectioncolor}{\textbf{Who is supporting/hosting/maintaining the dataset?
    }
    } \\
    %%%
    The authors of Lean Finder. \\
    %%%

    \textcolor{\sectioncolor}{\textbf{How can the owner/curator/manager of the dataset be contacted (e.g., email
    address)?
    }
    } \\
    %%%
    Contact information will be provided with the dataset release. \\
    %%%

    \textcolor{\sectioncolor}{\textbf{Is there an erratum?
    }
    If so, please provide a link or other access point.
    } \\
    %%%
    N/A \\
    %%%

    \textcolor{\sectioncolor}{\textbf{Will the dataset be updated (e.g., to correct labeling errors, add new
    instances, delete instances)?
    }
    If so, please describe how often, by whom, and how updates will be
    communicated to users (e.g., mailing list, GitHub)?
    } \\
    %%%
    Yes, we will frequently update the dataset to include new Lean statements in mathlib and other Lean related repositories we used.  \\
    %%%

    \textcolor{\sectioncolor}{\textbf{If the dataset relates to people, are there applicable limits on the
    retention of the data associated with the instances (e.g., were individuals
    in question told that their data would be retained for a fixed period of
    time and then deleted)?
    }
    If so, please describe these limits and explain how they will be enforced.
    } \\
    %%%
    N/A \\
    %%%

    \textcolor{\sectioncolor}{\textbf{Will older versions of the dataset continue to be
    supported/hosted/maintained?
    }
    If so, please describe how. If not, please describe how its obsolescence
    will be communicated to users.
    } \\
    %%%
    Yes, previous versions will remain accessible for reproducibility. \\
    %%%

    \textcolor{\sectioncolor}{\textbf{If others want to extend/augment/build on/contribute to the dataset, is
    there a mechanism for them to do so?
    }
    If so, please provide a description. Will these contributions be
    validated/verified? If so, please describe how. If not, why not? Is there a
    process for communicating/distributing these contributions to other users?
    If so, please provide a description.
    } \\
    %%%
    Yes, contributions will be encouraged through a collaborative platform (e.g., GitHub).\\
    %%%

    \textcolor{\sectioncolor}{\textbf{Any other comments?
    }} \\
    %%%
    The dataset’s maintainers are committed to ensuring its long-term usability and relevance for scientific research.
    %%%

\vspace{-2ex}
\section{Misc.}
\textbf{URL to our web service, model, and benchmark.} \url{https://leanfinder.github.io} We will continually update our system to support more Lean projects and a growing database of Lean statements.

\textbf{Author statement \& license information.} We the authors bear all responsibility in case of violation of rights.

\textbf{Dataset Structure.}
All files are stored in the JSONL format. Each input modality will be stored in a separate JSONL file. Each sample will contain the query and the ground truth Lean statements. We will also release our corpus of Lean statements.